\chapter*{\centerline{ЗАКЛЮЧЕНИЕ}}                      % Заголовок
\addcontentsline{toc}{chapter}{ЗАКЛЮЧЕНИЕ}  % добавляем его в оглавление

%% Согласно ГОСТ Р 7.0.11-2011:
%% 5.3.3 В заключении диссертации излагают итоги выполненного исследования, рекомендации, перспективы дальнейшей разработки темы.
%% 9.2.3 В заключении автореферата диссертации излагают итоги данного исследования, рекомендации и перспективы дальнейшей разработки темы.
%% Поэтому имеет смысл сделать эту часть общей и загрузить из одного файла в автореферат и в диссертацию:
В ходе диссертационного исследования были выполнены следующие задачи:

\begin{enumerate}[label=\arabic*)] % https://tex.stackexchange.com/a/476052/104425
\item Рассмотрены теоретические аспекты влияния эмпатии в контексте теории игр.
\item Разработан фреймворк для проведения экспериментов для оценки влияния эмоций в больших языковых моделях  на примере игр «Диктатор», «Дилемма узника» и «Общественное благо».
\item Проведены эксперименты, связанные с эмоциональным воздействием, в исследуемых играх.
\item Выявлены модели, наиболее соответствующие выравниванию с эмоциональным поведением человека.
\end{enumerate}

Разработанный фреймворк интегрирует эмоции в процесс принятия решений LLM (языковых моделей) в теории игр, позволяя тонко моделировать человекоподобное поведение искусственных агентов. Этот фреймворк учитывает эмоциональные факторы, влияющие на принятие решений, такие как страх, гнев, радость и грусть. Кроме того, он включает в себя элементы эмпатии, позволяя искусственным агентам лучше понимать и реагировать на эмоциональные состояния других участников игры. Благодаря этому, искусственные агенты могут демонстрировать более реалистичное, человекоподобное и социально адаптированное поведение в различных игровых сценариях и стратегических ситуациях.

В различных стратегических играх было доказано, что текущие языковые модели еще не готовы к непосредственному принятию решений из-за эмоциональных и стратегических предубеждений. Особенно сильно это влияние проявляется в открытых и небольших по размеру языковых моделях. В ходе исследования было выявлено, что эти модели склонны к иррациональным решениям, продиктованным эмоциями, а не логикой и стратегическим мышлением. Кроме того, они часто не способны проявлять эмпатию и учитывать эмоциональные состояния других участников игры, что приводит к неоптимальным результатам. Это ставит под сомнение их способность эффективно функционировать в сложных игровых ситуациях без дополнительной настройки и обучения.

Эксперименты показывают, что эмоциональный промптинг с использованием гнева может нарушить любое выравнивание, даже в GPT-4, демонстрируя человекоподобную эмоциональную реакцию. Это означает, что даже самые продвинутые языковые модели, такие как GPT-4, не застрахованы от влияния эмоций на их поведение. Более того, в этих ситуациях модель не проявляла эмпатии и не учитывала эмоциональное состояние других участников, что приводило к деструктивным последствиям. Это подчеркивает важность учета эмоционального фактора и развития эмпатии при разработке и использовании языковых моделей в различных приложениях.
